\section{Transfer to a topological counterexample}
\label{transfer:section}

Theorem~\ref{obs:nonmembership} completes the algebraic construction
on $\mathcal N$. We now turn that object into a topological
counterexample. The comparison of Section~\ref{comp:section} takes
place over the integral stack $\mathcal M$. Recall the coefficient maps
\[
 \mathcal N\xleftarrow{r}\mathcal M_k
       \xrightarrow{e}\mathcal M_0\xrightarrow{i}\mathcal M,
\]
where $r$ removes the central automorphisms, $e$ is the finite field
extension, and $i$ is the characteristic-$p$ inclusion. The main
algebraic step is to reverse a hypothetical finite construction through
these maps. The comparison triangles for reduction and removal of the
center may be nonsplit; their boundary maps will instead be nilpotent.

\Needspace{12\baselineskip}
\begin{proposition}[Transfer criterion]
\label{transfer:criterion}
Let $p>8$, with $D=2(p-1)>12$, and choose the comparison of
Theorem~\ref{comp:algebraicity}. Suppose $U\in\mathbf D(\mathcal N)$ satisfies
\begin{equation}
 \pi_iU\text{ is finite for every }i\in\Z,\qquad
 \widehat U\in\Thick(\widehat T'),\qquad
 U\notin\Thick(T').
 \label{transfer:conditions}
\end{equation}
Choose an $E(2)$-local spectrum $Z_{\rm sp}$ whose image under
$\mathcal F$ is
\[
 \widetilde U=i_*e_*r^*U\in\mathbf D(\mathcal M).
\]
Then
\[
 X_{\rm sp}=\tau_{\geq0}Z_{\rm sp}
\]
is a connective $p$-complete fp-spectrum of fp-type at most $2$,
and
\[
 L_{E(2)}X_{\rm sp}\simeq Z_{\rm sp},\qquad
 X_{\rm sp}\notin
 \Thick_{\Sp_p^\wedge}(b_p^\wedge).
\]
Here $b$ is the standard regular quotient of
Proposition~\ref{comp:standard-cohomology}.
\end{proposition}

The proposition is stated for any $U$ with these three properties;
the object already constructed will be its final application. The
proof has two parts. First we construct derived reduction
$Li^*:\mathbf D(\mathcal M)\to\mathbf D(\mathcal M_0)$ and,
writing $F=e_*r^*U$, establish the implications
\[
 \begin{split}
 Li^*i_*F\in\Thick(T_0)
 &\Longrightarrow F\in\Thick(T_0)\\
 &\Longrightarrow r^*U\in\Thick(T_k)
 \Longrightarrow U\in\Thick(T').
 \end{split}
\]
The last conclusion contradicts the third hypothesis. We then use the
comparison of the pure unit to prove that the connective cover is
$p$-complete and satisfies the \emph{total} finiteness test for fp-type.

\subsection{Derived reduction and its obstruction}

Call a sheaf, or a term of a complex, \emph{$p$-monic} if multiplication
by $p$ is injective on it. Derived reduction can be constructed using
such objects even though the characteristic-$p$ inclusion is not flat.
The following triangle records exactly how reducing an object already
in characteristic $p$ differs from returning that object unchanged.

\begin{lemma}[Reduction triangle]
\label{transfer:reduction}
There is an exact functor
\[
 Li^*:\mathbf D(\mathcal M)\longrightarrow
                 \mathbf D(\mathcal M_0)
\]
which on a termwise $p$-monic complex is ordinary reduction modulo $p$.
For $F\in\mathbf D(\mathcal M_0)$ it has a natural comparison triangle
\begin{equation}
 F[1]\longrightarrow Li^*i_*F\longrightarrow F
      \xrightarrow{\delta_F}F[2].
 \label{transfer:reduction-triangle}
\end{equation}
\end{lemma}

\begin{proof}
We first construct enough $p$-monic epimorphisms of sheaves. On the
whole graded $\BP$ stack, the cofree-on-the-base comodule $\Gamma_P$
is a union of graded dualizable subcomodules which are finite free over
the left base. Specifically, take the left spans of the arrow monomials
of bounded weight. There are finitely many such monomials at each
bound. The coproduct preserves these spans because it is homogeneous,
all base and arrow coefficients have nonnegative weight, and left
scalars remain in the outer factor of the coaction.

For a graded comodule $M$, the diagonal tensor $\Gamma_P\otimes_PM$
is cofree on the underlying graded module of $M$. This is also the
projection formula for the affine chart
$[\Spec P/\Gm]\to\mathcal M_{\rm all}$ of the whole stack.
A homogeneous element of $M$ therefore corresponds to a map from a
grading shift of the unit into this diagonal tensor. It factors through
$N\otimes_PM$ for one of the finite-free subcomodules $N$:
primitives commute with the filtered union, or equivalently the given
coaction expression has finite support. The inclusions of these spans
are split on underlying left modules, so tensoring with $M$ preserves
the injections and the factorization respects coactions. Duality turns
it into a map from a shifted $N^\vee$ to $M$. Evaluation at the
counit restricted to $N$ shows that the original element is in its
image. Summing these maps gives an epimorphism from a sum of
dualizables to every graded comodule.

To apply this on $\mathcal M$, extend a quasi-coherent sheaf by
pushforward from the height-open into the whole stack. This
pushforward is quasi-coherent: on the flat affine charts the open is
quasi-compact, and its pushforward commutes with flat base change.
Restrict the epimorphism just constructed. The restricted dualizables
are still $p$-monic, so $\mathcal M$ has enough $p$-monic
epimorphisms.

For a twisted periodic complex $C$, use disk sums on these generators
to choose a functorial degreewise epimorphism $P_C\to C$. Its
kernel $K_C$ is termwise $p$-monic, as is $P_C$. The natural map
\[
 \cofib(K_C\longrightarrow P_C)\longrightarrow C
\]
is a quasi-isomorphism, and its source is termwise $p$-monic.
The same construction works within the full subcategory of termwise
$p$-monic complexes, so that subcategory has the same
quasi-isomorphism localization as all complexes. An acyclic complex
in the subcategory stays acyclic after reduction: apply the short
exact multiplication-by-$p$ sequence. Since the subcategory is
closed under shifts and cones, ordinary reduction on it induces an
exact functor on its localized homotopy category. This constructs
$Li^*$.

For $C=i_*F$, the snake sequence for multiplication by $p$ is
\[
 0\longrightarrow F\longrightarrow K_C/p
  \longrightarrow P_C/p\longrightarrow F\longrightarrow0
\]
as a sequence of complexes over $\mathcal M_0$. The cone on the
middle map is a representative of $Li^*i_*F$. It extends the
cokernel $F$ by the shifted kernel $F[1]$, giving
\eqref{transfer:reduction-triangle}. The iterated boundary is natural
for chain maps and hence descends naturally to the localization.
\end{proof}

This supplies the first boundary map whose nilpotence will reverse a
coefficient change. The other comparison triangle was constructed in
Proposition~\ref{geom:relative-pushforward}:
\begin{equation}
 U\longrightarrow Qr^*U\longrightarrow U[-1]
       \longrightarrow U[1],\qquad Q=Rr_*.
 \label{transfer:center-triangle}
\end{equation}
After a shift, its last map is a natural transformation
$\epsilon_U:U\to U[2]$. We will use $\delta_F$ for reduction and
$\epsilon_U$ for removal of the center. The next lemma explains why
nilpotence of these transformations suffices to recover thick
membership from the middle terms.

\begin{lemma}[Nilpotent comparison obstructions]
\label{transfer:nilpotence}
Let $\delta:\id\to[2]$ be a natural transformation on a triangulated
category. For an object $A$, consider the property that there is an
integer $N\geq1$ such that
\[
 \delta_{A[j]}^N:A[j]\longrightarrow A[j+2N]
 \quad\text{is zero for every }j\in\Z,
\]
where the power denotes the successive natural transformations on
the intervening shifts. No compatibility between $\delta_{A[j]}$ and
$\delta_A[j]$ is assumed. This property is thick.

Suppose in addition that the cones of the maps $\delta_{A[j]}$ are
shifts of $B(A[j])$, where $B$ is an exact functor. If $A$ has the
property, then $A\in\Thick(B(A))$.
\end{lemma}

\begin{proof}
Shifts and retracts preserve the stated property. Consider a triangle
$A\to B_1\to C$ and exponents $a,c$ which work on every shift
of $A,C$, respectively. Naturality makes the composite of the
$c$-fold power on $B_1[j]$ with the map to $C[j+2c]$ vanish.
Exactness of represented sequences factors that power through
$A[j+2c]$. Another $a$ applications of $\delta$ kill the
factorization, again by naturality. Thus $a+c$ works uniformly
on all shifts of $B_1$.

For the second assertion, the cone of each factor in
$\delta_A^N$ belongs to $\Thick(B(A))$: exactness identifies
$B(A[j])$ with a shift of $B(A)$. Repeated octahedra show that
the cone of their composite also belongs there. If the composite
is zero, that cone is $A[2N]\oplus A[1]$. Taking a retract and
then a shift proves the claim.
\end{proof}

For \eqref{transfer:reduction-triangle} the cone of $\delta_F$ is
$(Li^*i_*F)[1]$. For the shifted boundary $\epsilon_U$ of
\eqref{transfer:center-triangle}, it is $(Qr^*U)[2]$.
Both middle-object functors are exact, so the lemma applies to each.

\subsection{Nilpotence on the open and complete pieces}

We next prove nilpotence for the particular objects in the transfer
criterion. Fracture reduces this to the $t$-open and the completion.
On an object pushed forward from the $t$-open, both obstructions
vanish. Indeed Lemma~\ref{complete:open-splitting} gives a finite
decomposition into pure period-residue summands. For $\mathcal N$
these are vector spaces; for $\mathcal M_0$ they are smooth
representations of the procyclic center. Each pure summand has no
self-map to its $2$-shift. Naturality against the inclusions of the
summands therefore kills the obstruction on their direct sum, and
the same reasoning applies to every shift. The vector spaces or
representations themselves need not have finite dimension: there are
only finitely many period residues.

By Corollary~\ref{complete:degree-two}, the complete generators
$\widehat T'$, $\widehat T_k$, and $\widehat T_0$ likewise
have no self-maps to their $2$-shifts, and the same holds after every
shift. Hence the appropriate obstruction
is nilpotent, uniformly over all shifts, on each of their thick
subcategories. To use this for the reduction comparison we must also
control the completed pullback through the center.

\begin{lemma}[Completed central pullback]
\label{transfer:central-completion}
Suppose $U\in\mathbf D(\mathcal N)$ and
$\widehat U\in\Thick(\widehat T')$. Set
$F=e_*r^*U\in\mathbf D(\mathcal M_0)$. Then
\[
 \widehat F\in\Thick(\widehat T_0).
\]
The comparison obstructions of
\eqref{transfer:center-triangle} and
\eqref{transfer:reduction-triangle} are uniformly nilpotent over all shifts of
$U$ and $F$, respectively.
\end{lemma}

\begin{proof}
The fiber of $U\to\widehat U$ is $t$-invertible by
Proposition~\ref{complete:fracture}. Pullback preserves this
property, and completion kills such objects. Therefore
\[
 \widehat{r^*U}\simeq\widehat{r^*(\widehat U)}.
\]
The exact functor obtained by composing $r^*$ with completion
reduces the assertion to the generator. Its value on
$\widehat T'$ is its value on $T'$, by the same invertible-fiber
argument. The finite reductions are $r^*T'_d$, where
$T'_d=T'/t^d$.

At each level the full coinduction description of
Proposition~\ref{geom:coinduction} gives an exact sequence
\begin{equation}
 0\longrightarrow r^*T'_d\longrightarrow T_{k,d}
       \xrightarrow{c-1}T_{k,d}\longrightarrow0,
 \qquad T_{k,d}=T_k/t^d.
 \label{transfer:central-difference}
\end{equation}
Here $c$ is a chosen topological generator of the center
$Z=1+p\Z_p$.
We include the surjectivity check, which uses characteristic $p$.
A locally constant function on $G_{\rm up}/\Lambda$ factors
through a finite quotient. Refine this quotient until each central
period has increased by a multiple of $p$. This is possible because
the center embeds in the profinite group and its central $p$-power
images meet the tame $\Lambda$-images trivially. On each refined
central cycle, the sum of the repeated values is zero in
characteristic $p$. The function can therefore be integrated as a
discrete difference. This proves surjectivity of $c-1$; its kernel
is the rigidified coinduction object.

Take homotopy limits of \eqref{transfer:central-difference}. The
result is a triangle expressing $\widehat{r^*T'}$ using two copies
of $\widehat T_k$. Thus
$\widehat{r^*U}\in\Thick(\widehat T_k)$. This argument uses
compatible finite reductions, not commutation of $r^*$ with
arbitrary homotopy limits. Restriction along the finite field extension
commutes with the limits and sends $\widehat T_k$ to a finite
sum of copies of $\widehat T_0$. Hence it gives the claimed
membership for $\widehat F$.

We have already proved nilpotence on the open pieces and on the thick
subcategories of the complete generators. Apply it to the fracture
triangles
\[
 U\longrightarrow\widehat U\oplus LU
               \longrightarrow L\widehat U,
 \qquad
 F\longrightarrow\widehat F\oplus LF
               \longrightarrow L\widehat F.
\]
The open terms have zero obstruction, and the complete terms have
nilpotent obstruction. Lemma~\ref{transfer:nilpotence} supplies a
uniform exponent over all shifts of $U$ and $F$.
\end{proof}

The completed central pullback is now controlled by the original
completed chart, and both comparison boundaries have uniform nilpotence
exponents. We can therefore reverse the coefficient changes at the
level of thick membership.

\begin{lemma}[Excluding the reduced construction]
\label{transfer:reduced-exclusion}
If $U$ satisfies \eqref{transfer:conditions}, then
\[
 Li^*(i_*e_*r^*U)\notin\Thick(T_0).
\]
\end{lemma}

\begin{proof}
Write $F=e_*r^*U$. If the displayed membership held, the
reduction comparison triangle, Lemma~\ref{transfer:central-completion},
and Lemma~\ref{transfer:nilpotence} would give
$F\in\Thick(T_0)$. Exact scalar extension then gives
\[
 e^*F\in\Thick(e^*T_0)=\Thick(T_k).
\]
The object $r^*U$ is a retract of $e^*e_*r^*U$ by separability
of the finite field extension. Hence $r^*U\in\Thick(T_k)$.
Apply the exact functor $Q$ and use
$QT_k=T'$ from Proposition~\ref{geom:relative-pushforward}.
It follows that $Qr^*U\in\Thick(T')$. Nilpotence of the
central comparison obstruction and
Lemma~\ref{transfer:nilpotence} now imply
$U\in\Thick(T')$, contrary to \eqref{transfer:conditions}.
\end{proof}

\subsection{The connective cover and the finite test}

\begin{proof}[Proof of Proposition~\ref{transfer:criterion}]
Choose $Z_{\rm sp}$ as in the statement, using essential
surjectivity of Theorem~\ref{comp:algebraicity}. By
Proposition~\ref{comp:generators}, its ordinary homotopy groups
identify with sections homotopy of $\widetilde U$:
\[
 \begin{aligned}
 \pi_i Z_{\rm sp}
 &\cong[\Sigma^iL_ES^0,Z_{\rm sp}]_{h\Sp_E}\\
 &\cong[\mathcal O[i],i_*e_*r^*U]_{\mathbf D(\mathcal M)}\\
 &\cong\pi_i Qr^*U.
 \end{aligned}
\]
These are identifications of abelian groups, since a triangulated
functor is additive. The last identification uses the acyclic
resolution and closed-base
change of Section~\ref{geom:section}. For example, $i_*e_*$ sends
termwise injectives to complexes of sections-acyclic objects, and the
finite sections bounds justify the computation on periodic complexes.
The comparison triangle \eqref{transfer:center-triangle} and the
first condition in \eqref{transfer:conditions} show that these groups
are finite $p$-groups in every degree.

Set $X_{\rm sp}=\tau_{\geq0}Z_{\rm sp}$. It is connective,
and its homotopy groups are finite $p$-groups. Its finite Postnikov
stages are $p$-complete, and $p$-complete spectra are closed under
limits. The Postnikov limit consequently proves that $X_{\rm sp}$
is $p$-complete.

Next, $H\F_p$ is $E(2)$-acyclic. One can see this directly in the
same module calculation as Proposition~\ref{comp:bp-comodule}:
connective ring truncation gives its $\BP$-module structure and
\[
 \BP_*H\F_p=\Gamma_P\otimes_{\eta_R,P}\F_p.
\]
The left positive coefficient actions vanish successively, using
$\eta_Rv_i\equiv v_i\pmod{(p,v_1,\ldots,v_{i-1})}$.
Landweber base change to the chart with $v_2$ invertible therefore
gives zero. The negative tail $\tau_{<0}Z_{\rm sp}$ is a filtered
colimit of finite Postnikov ranges with finite $p$-group homotopy.
A composition series of each finite $p$-group expresses its
Eilenberg--Mac Lane spectrum as a finite construction from $H\F_p$.
Consequently each range is a finite construction from shifts of
$H\F_p$ and is $E(2)$-acyclic. Thus the whole tail is acyclic and
\[
 L_{E(2)}X_{\rm sp}\simeq Z_{\rm sp}.
\]

Let $V$ be any finite type-three spectrum. Such spectra exist by
the periodicity theorem, taking successive self-map cofibers
\cite[Theorem~1.5.4 and Corollary~2.3.5]{Ravenel1992}. The chromatic
Bousfield-class calculation \cite[Theorem~7.3.2(d)]{Ravenel1992}
\[
 \langle E(2)\rangle
   =\langle K(0)\vee K(1)\vee K(2)\rangle
\]
shows that both $V$ and its Spanier--Whitehead dual are
$E(2)$-acyclic. Indeed finite duality over each Morava $K$-theory
field preserves the vanishing that defines type three. If $DV$ is the
Spanier--Whitehead dual and $W$ is any spectrum, finite duality gives
\[
 [W,V\wedge Z_{\rm sp}]\cong[DV\wedge W,Z_{\rm sp}]=0,
\]
because $DV\wedge W$ is $E(2)$-acyclic and $Z_{\rm sp}$ is local. Hence
$V\wedge Z_{\rm sp}=0$. Consequently
\[
 V\wedge X_{\rm sp}
   \simeq\Sigma^{-1}\bigl(V\wedge\tau_{<0}Z_{\rm sp}\bigr).
\]
The left side is bounded below and has finite homotopy in each
degree, since $X_{\rm sp}$ is connective with finite homotopy and
$V$ has finitely many cells. The right side is bounded above,
since its tail is bounded above and $V$ has finitely many cells.
It follows that $\pi_*(V\wedge X_{\rm sp})$ is finite in total.
Since $X_{\rm sp}$ is connective and $p$-complete,
Mahowald--Rezk's criterion \cite[Proposition~3.2]{MR1999} now
proves fp-ness as well as fp-type at most $2$; fp-ness is not an
additional assumption in that criterion. This connective-cover step
is also related to their general formulation in Proposition~3.5 and
the chromatic construction in Section~7.2(2).

Finally suppose that $X_{\rm sp}$ belonged to
$\Thick_{\Sp_p^\wedge}(b_p^\wedge)$. The full category of
$p$-complete spectra is stable, so its finite constructions remain
finite constructions of spectra. Exact localization gives
\[
 Z_{\rm sp}\in\Thick(L_Eb_p^\wedge).
\]
Apply the triangulated equivalence $\mathcal F$. It preserves shifts,
distinguished triangles, and retracts, while
Proposition~\ref{comp:generators} identifies the generator with
$T_{\Z_p}$. Hence
\[
 \widetilde U\in\Thick(T_{\Z_p}).
\]
The pure sheaf $T_{\Z_p}$ is $p$-monic and has reduction $T_0$ by
Proposition~\ref{comp:bp-comodule}. Exactness of $Li^*$ therefore gives
$Li^*\widetilde U\in\Thick(T_0)$. This contradicts
Lemma~\ref{transfer:reduced-exclusion} and completes the proof.
\end{proof}

For the glued object, Proposition~\ref{glue:object} supplies the first
two conditions in \eqref{transfer:conditions}, and
Theorem~\ref{obs:nonmembership} supplies the third. The transfer
criterion therefore produces the spectrum used in
Theorem~\ref{thm:main}. We next establish its exact fp-type and its
localization properties.
